\chapter{3D Gaussian Splatting}
\label{ch:theory2}

Neural radiance fields (NeRF) showed that a scene can be encoded as a continuous volumetric function approximated by a neural network, enabling high-quality view synthesis from posed images. Training and rendering were initially slow. \textbf{3D Gaussian Splatting} (3DGS), introduced by Kerbl et al.\ \cite{kerbl20233dgs}, replaces the neural field by a set of \textbf{explicit anisotropic Gaussians} in space. Kerbl et al.\ show that this representation delivers \textbf{high visual quality together with real-time rendering}: their method matches or exceeds prior radiance-field quality while running at interactive frame rates on desktop GPUs \cite{kerbl20233dgs}. In this thesis, 3DGS is used as an \textbf{imported} representation (\texttt{.ply} files) rendered with a custom Metal pipeline for \textbf{comparison and viewing}; mesh-based USDZ assets remain the medium for interactive furniture editing and for Blender export. Figure~\ref{fig:ch3_nerf_pipeline} illustrates the NeRF rendering formulation that motivated later explicit real-time representations such as 3DGS.

\begin{figure}[H]
\centering
\thesisfig{0.9\textwidth}{figures/fig3_1_nerf_overview.pdf}
\caption{NeRF overview: sample points along camera rays are fed to a neural MLP predicting colour and density; volume rendering accumulates pixel colours; a photometric loss drives backpropagation (adapted from \cite{mildenhall2020nerf}).}
\label{fig:ch3_nerf_pipeline}
\end{figure}

\section{Representation: 3D Gaussian Primitives}
\label{sec:3dgs_representation}

Each primitive stores a mean $\boldsymbol{\mu} \in \mathbb{R}^3$, a covariance $\boldsymbol{\Sigma}$ describing elongation and orientation, an opacity $\alpha$, and colour parameters. Following Kerbl et al.\ \cite{kerbl20233dgs}, $\boldsymbol{\Sigma}$ is built from a rotation matrix $\mathbf{R}$ (quaternion $\mathbf{q}$) and diagonal scale $\mathbf{S}=\operatorname{diag}(\mathbf{s})$ with positive entries: $\boldsymbol{\Sigma} = \mathbf{R}\mathbf{S}\mathbf{S}^\top\mathbf{R}^\top$. Colour may be view-dependent via spherical harmonics; the implementation described later also accepts per-Gaussian RGB in linear space.

The spatial falloff is
\[
G(\mathbf{x}) = \exp\bigl(-\tfrac{1}{2}(\mathbf{x}-\boldsymbol{\mu})^\top \boldsymbol{\Sigma}^{-1} (\mathbf{x}-\boldsymbol{\mu})\bigr).
\]

\begin{figure}[H]
\centering
\thesisfig{\textwidth}{figures/fig3_2_gaussian_splatting_projection.pdf}
\caption{Conceptual illustration of 3D Gaussian projection and splatting: each 3D Gaussian is projected onto the image plane; footprints are depth-sorted and alpha-composited \cite{kerbl20233dgs}.}
\label{fig:gaussian_project}
\end{figure}

\section{Projection, Rasterisation, and Real-Time Rendering}
\label{sec:3dgs_splatting}

Rendering splats each Gaussian onto the image \cite{kerbl20233dgs}:

\begin{enumerate}
    \item \textbf{Project} $\boldsymbol{\mu}$ with the camera matrix.
    \item \textbf{Project covariance} to $2\times2$ image-space covariance using the Jacobian $\mathbf{J}$ of the perspective mapping at $\boldsymbol{\mu}$.
    \item \textbf{Sort} primitives (globally or per tile) for approximate depth order.
    \item \textbf{Rasterise} with alpha blending, accumulating colour and transmittance along rays.
\end{enumerate}

At each pixel, ordered compositing follows
\[
\mathbf{c}_{\mathrm{out}} = \alpha \cdot \mathbf{c}_{\mathrm{G}} + (1-\alpha) \mathbf{c}_{\mathrm{in}}.
\]

\section{Differentiable Rendering and Optimisation}
\label{sec:3dgs_differentiable}

To fit Gaussians to photographs, the renderer must be \textbf{differentiable}. Gradients of a photometric loss propagate backwards through alpha compositing, rasterisation, and projection, updating $\boldsymbol{\mu}$, scales, and rotations. Optimisation typically alternates gradient steps with \textbf{adaptive density control} (splitting large Gaussians, culling transparent ones) \cite{kerbl20233dgs}. Initialisation often starts from a sparse point cloud (e.g.\ from COLMAP SfM), linking 3DGS training back to the photogrammetry chain in Chapter~\ref{ch:theory1}. Full training remains \textbf{off-device} in this application; the iOS code imports and renders only the optimised \texttt{.ply} result.

\section{Meshes versus Gaussians: Editability and Role in the Application}
\label{sec:3dgs_vs_mesh}

\textbf{Meshes (USDZ from Object Capture)} expose \textbf{explicit connectivity}: faces can be simplified or exported; axis-aligned bounding boxes align with RoomPlan furniture detections; RealityKit loads USDZ natively. \textbf{3D Gaussians} \cite{kerbl20233dgs} do not define a thin surface: coherent editing requires transforming thousands of primitives; boolean operations and CAD-style placement are ill-defined without mesh conversion.

The product decision in this thesis is therefore:

\begin{itemize}
    \item \textbf{Authoring and replacement}: USDZ meshes only (Asset~A).
    \item \textbf{Visualisation and comparison}: optional \texttt{.ply} splats via Metal (Assets~B and~C).
    \item \textbf{Desktop refinement}: Blender export carries glTF room structure and USDZ furniture copies; splat \texttt{.ply} files are not bundled in the ZIP and remain a separate in-app benchmark unless imported manually in Blender.
    \item \textbf{Training}: off-device for Asset~B; feedforward inference on a LAN Mac for Asset~C.
\end{itemize}

\begin{table}[htbp]
\centering
\caption{Mesh-based Object Capture output versus imported 3D Gaussian Splatting in this application.}
\label{tab:mesh_vs_3dgs}
\begin{tabular}{|p{0.22\textwidth}|p{0.34\textwidth}|p{0.34\textwidth}|}
\hline
\textbf{Criterion} & \textbf{Mesh (USDZ)} & \textbf{3DGS (\texttt{.ply})} \\
\hline
Geometry & Explicit triangles & Cloud of Gaussians \\
\hline
Rendering & RealityKit / Metal mesh & Splatting \cite{kerbl20233dgs} \\
\hline
RoomPlan alignment & Scale-to-bounds, grounding & Manual alignment if needed \\
\hline
Editing in app & Yes (replace, transform) & View / compare only \\
\hline
Blender export & Included (copies) & Not included in export ZIP \\
\hline
\end{tabular}
\end{table}

\section{Three Object-Level Reconstruction Paths}
\label{sec:three_assets}

Experimental comparisons in this thesis evaluate physical objects or room regions through three pipelines with different input labour, compute location, and editability. Only Asset~A integrates with interactive furniture replacement; Assets~B and~C share the Metal splat viewer.

\textbf{Asset~A (Object Capture).} Guided multi-view orbit on device; photogrammetry session produces metric textured USDZ for the furniture library and Blender export copies.

\textbf{Asset~B (offline 3DGS).} Fifty to two hundred or more photographs with poses from COLMAP-class SfM on a workstation; GPU training yields a standard \texttt{.ply} imported into the app. Highest plenoptic quality when capture discipline is good; relative scale unless metric constraints are added.

\textbf{Asset~C (SHARP).} Single JPEG posted to a user-controlled Mac on the local network; feedforward model returns a pruned \texttt{.ply} ($\leq 300\,000$ Gaussians) for mobile viewing. Minimises capture friction at the cost of single-view ambiguity and network dependency.

\begin{table}[htbp]
\centering
\caption{Three reconstruction paths compared (object-level focus).}
\label{tab:three_assets}
\small
\begin{tabular}{|p{2.1cm}|p{3.1cm}|p{3.1cm}|p{3.1cm}|}
\hline
\textbf{Criterion} & \textbf{A: Object Capture} & \textbf{B: Offline 3DGS} & \textbf{C: SHARP} \\
\hline
Input & Multi-view orbit & Multi-view + SfM poses & Single image \\
\hline
Where it runs & iPhone/iPad & External GPU & Mac on LAN \\
\hline
Time-to-result & Minutes & Minutes--hours & $\approx 10\,$s (warm server) \\
\hline
Output & USDZ mesh & \texttt{.ply} splat & \texttt{.ply} splat \\
\hline
Edit in app & Yes & View only & View only \\
\hline
Renderer & RealityKit & Metal splats & Metal splats \\
\hline
\end{tabular}
\end{table}

\section{Limitations of Classical 3DGS for Mobile Capture}
\label{sec:3dgs_limitations}

Despite its visual quality and real-time rendering speed, the classical 3DGS pipeline~\cite{kerbl20233dgs} carries practical constraints for mobile, user-initiated capture: multi-view input, offline SfM, per-scene optimisation, and typically relative scale from COLMAP initialisation. These constraints motivate feedforward SHARP (Asset~C) for rapid previews while keeping full training off-device for Asset~B.

\section{SHARP: Feedforward Single-Image 3D Gaussian Splatting}
\label{sec:sharp}

\textbf{SHARP} (\emph{Sharp Monocular View Synthesis})~\cite{sharp-monocular-view} predicts a complete Gaussian field from one RGB image in a single forward pass, without per-scene optimisation. It builds on a monocular metric depth prior, initialises on the order of $10^6$ Gaussians, and decodes attributes compatible with Kerbl et al.'s rasterisation model~\cite{kerbl20233dgs}. Table~\ref{tab:3dgs_vs_sharp} contrasts offline optimised 3DGS (Asset~B) with SHARP (Asset~C).

\begin{table}[h]
\centering
\caption{Classical 3DGS~\cite{kerbl20233dgs} versus SHARP~\cite{sharp-monocular-view}.}
\label{tab:3dgs_vs_sharp}
\setlength{\tabcolsep}{4pt}
\begin{tabular}{p{2.8cm} p{4.3cm} p{4.3cm}}
\toprule
\textbf{Aspect} & \textbf{Classical 3DGS (Asset~B)} & \textbf{SHARP (Asset~C)} \\
\midrule
Input             & Multi-view images + poses & Single image \\
Pose estimation   & COLMAP SfM (offline)      & Not required \\
Inference mode    & Per-scene optimisation    & Single forward pass \\
Inference time    & Minutes to hours          & ${<}$1\,s (GPU) / ${\approx}$10\,s (Apple Silicon) \\
Mobile suitability & Import \texttt{.ply} only & LAN server $\rightarrow$ import \texttt{.ply} \\
\bottomrule
\end{tabular}
\end{table}

\subsection{Integration in This Application}
\label{subsec:sharp_integration}

The application integrates SHARP using a client--server architecture: a Flask server (\texttt{sharp\_server.py}) on the local network exposes \texttt{/sharp/infer}; the iOS client (\texttt{SHARPService}, \texttt{SHARPCaptureView}) POSTs a JPEG with extended timeouts and displays the returned pruned \texttt{.ply} in the Metal viewer. Implementation details appear in Chapter~\ref{ch:implementation}.

This positions the application to demonstrate a reconstruction spectrum: \textbf{RoomPlan} (parametric, real-time) $\rightarrow$ \textbf{Object Capture} (photogrammetric mesh, on-device) $\rightarrow$ \textbf{offline 3DGS} (Asset~B) $\rightarrow$ \textbf{SHARP} (Asset~C) $\rightarrow$ optional \textbf{Blender export} of the staged mesh-based room (Chapter~\ref{ch:implementation}).
